\chapter{Faults}\label{ch:faults}
\begin{goals}
\textbf{In this chapter} you open fault-slip data, draw it with slip arrows, P and T axes and right
dihedra, and read the stress tensor wuffle inverts from it.
\end{goals}

A fault measurement is a plane and a line on it: the fault plane, the slickenline (striation) on it, and,
when you know it, the sense in which the hanging wall moved. wuffle keeps the three together as one row of
a \term[fault dataset]{fault dataset}, draws the slip on the net, and inverts a set of faults for the stress
that would drive them.

\section{Fault data}\label{sec:fault-data}

A file reads as faults when it has the columns of a plane \emph{and} the columns of a line
(\cref{ch:data}):

\begin{center}\small
\begin{tabularx}{\linewidth}{@{}>{\raggedright}p{0.24\linewidth}L@{}}\toprule
\sffamily\bfseries holds & \sffamily\bfseries column names wuffle recognises (case, spaces and punctuation ignored) \\\midrule
the plane & \col{dip direction} (\col{dipdir}, \col{dd}, \col{dip azimuth}…) and \col{dip}; or \col{strike} and \col{dip}, read by the right-hand rule \\
the slickenline & \col{trend} (\col{tr}, \col{azimuth}, \col{bearing}…) and \col{plunge} (\col{pl}, \col{plg}…) \\
the sense & \col{sense}, \col{slip sense}, \col{kinematics}, \col{kin} or \col{movement} (optional) \\
\bottomrule\end{tabularx}
\end{center}

The sense column takes words or the one-character codes common in fault-slip files:

\begin{center}\small
\begin{tabularx}{\linewidth}{@{}>{\raggedright}p{0.2\linewidth}L@{}}\toprule
\sffamily\bfseries sense & \sffamily\bfseries written as \\\midrule
normal & \code{normal}, \code{N}, \code{2}, \code{-} \\
reverse & \code{reverse}, \code{inverse}, \code{thrust}, \code{R}, \code{I}, \code{1}, \code{+} \\
dextral & \code{dextral}, \code{right-lateral}, \code{D}, \code{3} \\
sinistral & \code{sinistral}, \code{left-lateral}, \code{S}, \code{L}, \code{4} \\
unknown & anything else, or empty: \code{unknown}, \code{?} \\
\bottomrule\end{tabularx}
\end{center}

Case doesn't matter. A row whose sense is unknown is still a fault: it is drawn, and it counts in the
statistics of its planes, but it carries no arrow and stays out of everything that needs the sense (P and
T axes, right dihedra, the stress inversion).

\begin{tip}
The trend and plunge only need to name the slickenline; the sense says which way along it the hanging
wall moved. Write the line as you measured it, lower-hemisphere as always, and let the sense column carry
the direction.
\end{tip}

\section{What a fault layer draws}\label{sec:fault-elements}

A fault layer opens in the layers tree into five elements (\cref{ch:elements}). A new one shows the first
two:

\begin{description}
\item[great circles] the fault planes;
\item[slip lines] the slickenline of each fault as a point, with an arrow for the hanging wall's motion;
\item[P and T axes] the shortening and extension axes of each fault;
\item[right dihedra] the boundary of the compressional and extensional fields the faults agree on;
\item[labels] a column's values beside each fault.
\end{description}

\shot{faults-net}{Twelve normal faults: their planes, the slickenlines with arrows pointing away from the
centre (the hanging walls went down), and the P (hollow) and T (filled) axes. The fault whose sense is
\code{?} has a point and no arrow, and no P or T.}

\subsection{Slip arrows}\label{sec:slip-arrows}

Each slickenline is drawn where its lower-hemisphere end plots, and the arrow on it follows
\term[Angelier's convention]{Angelier's convention}: it shows the motion of the hanging wall.

\begin{itemize}
\item An arrow pointing \emph{away from the centre} is a hanging wall that went down: a normal component.
\item An arrow pointing \emph{toward the centre} is a hanging wall that went up: a reverse component.
\end{itemize}

% BUG: stereonet core.js slipArrow only moves along the plunge, so a horizontal slickenline (plunge 0)
% gets a zero-length shaft and no visible arrow, and the strike-slip sense of a shallow slickenline is
% never drawn (a dextral and a sinistral fault on the same line look the same). The sense IS used for
% P/T, dihedra and the inversion. Book written to what happens.
The arrow shows only that up-or-down component. A purely strike-slip fault (a horizontal slickenline) gets
no arrow, and dextral and sinistral faults with the same slickenline look the same on the net; their sense
still counts for the P and T axes, the dihedra and the inversion. A vertical slickenline has no direction
on the net either, so it gets no arrow.

\subsection{P and T axes}

Each fault with a known sense gives a \term{P axis} and a \term{T axis}: the directions at 45° to the
fault plane, in the plane that holds the slickenline and the pole, on the shortening and the extension
side. The P axis is drawn hollow, the T axis filled. For a set of compatible faults the P axes cluster
round the direction of greatest compression and the T axes round the least.

\subsection{Right dihedra}\label{sec:dihedra}

\shot{faults-dihedra}{The same faults with \ui{right dihedra} on: the line where the faults' compressional
and extensional quadrants balance. Here it rings the centre, round the vertical shortening of these normal
faults.}

Each fault divides the sphere into four quadrants, two compressional and two extensional, by its plane
and the plane normal to its slip. The \term{right dihedra} element adds the faults' quadrants up over the
whole sphere and draws the line where compression and extension balance. The side holding the P axes is
where the most compressive stress can lie for every fault at once; switch \ui{P and T axes} on with it to
see which side that is.

\begin{caution}
On a rotated net (\cref{ch:rotation}) the right dihedra are not drawn: the field is computed in the
net's own unrotated frame. Turn the rotation off to see them.
\end{caution}

\section{The stress inversion}\label{sec:stress}

With four or more faults of known sense, the statistics card of a fault layer adds two lines from
the inversion of \textcite{michael1984}, which finds the stress tensor whose resolved shear direction on each
fault plane best matches its slip, assuming the shear stress has about the same magnitude on every fault:

\begin{description}
\item[\ui{σ1 σ2 σ3}] the principal stresses as lines: σ1 the most compressive, σ3 the least;
\item[\ui{φ}] the stress ratio $\phi = (\sigma_2 - \sigma_3)/(\sigma_1 - \sigma_3)$, from 0 (σ2 = σ3) to 1
(σ2 = σ1).
\end{description}

Hovering over \ui{σ1 σ2 σ3} tells you how many faults went in. With fewer than four of known sense, the
card says \ui{needs ≥ 4 faults with a known sense}.

\shot[0.7\linewidth]{faults-stats}{The statistics card of the twelve faults: the axial mean and eigen analysis
of their poles, then σ1 vertical and the stress ratio. The fault with sense \code{?} counts in the first
lines and not in the last two.}

The lines above them are those of any planes dataset (\cref{ch:stats}): the axial Fisher mean of the
fault poles, the eigen analysis, Woodcock's K and C \parencite{woodcock1977} and Vollmer's P, G and R
\parencite{vollmer1990}. The buttons under the card
work as for planes, so \ui{mean → dataset} keeps the mean fault plane as a dataset.

\begin{note}
The console's \code{stats} prints the n, the mean and the eigenvalues of a fault layer, not the stress;
the stress is on the card.
\end{note}

\begin{tryit}
\begin{enumerate}
\item Save these lines as \file{faults.csv}:
\begin{lstlisting}
dip direction,dip,trend,plunge,sense
000,60,000,60,normal
030,60,030,60,normal
060,60,060,60,normal
090,55,090,55,normal
120,60,120,60,normal
150,65,150,65,normal
180,60,180,60,normal
210,60,210,60,normal
240,50,240,50,normal
270,60,270,60,normal
300,60,300,60,normal
330,60,330,60,?
\end{lstlisting}
\item Open it (\menu{File>Open data or a project…}). It reads as faults: twelve great circles and their
slickenlines, eleven with arrows pointing outward.
\item In the tree, open the \ui{faults} layer and tick \ui{P and T axes}. The P axes gather near the
centre, the T axes near the rim: vertical shortening, horizontal extension.
\item Open the \ui{statistics} tab. σ1 is vertical (plunge 90) and $\phi$ is close to 0: a radial,
uniaxial extension, as twelve faults dipping away from one point should give.
\item Tick \ui{right dihedra}. A ring appears round the centre, with the P axes inside it.
\end{enumerate}
\end{tryit}
